%===================================================================================================
\section{Derived Requirements and Design Implications} %++++++++++++++++++++++++++++++++++++++++++++
  \label{Derived Requirements and Design Implications} %++++++++++++++++++++++++++++++++++++++++++++
%===================================================================================================


\begin{justify}
	Based on the issues above, this section derives requirements for novice-oriented test failure reporting.
	Each requirement is phrased as a constraint on the \emph{user-facing report}, independent of any specific implementation technique.
	The subsequent design chapter then maps these requirements onto concrete mechanisms (e.g., expression capture, value extraction, and framework wrappers).
\end{justify}


%===============================================================================
\subsection*{Requirements}
%===============================================================================


\begin{justify}
	\begin{description}
		\item[\textbf{R1: Stable summary structure.}]
			Every failing test should provide a compact, consistently ordered summary containing at least: (i) the tested expression (or an equivalent description of the behaviour under test), (ii) the expected value/condition, and (iii) the actual value/observed result.

		\item[\textbf{R2: Explicit tested expression.}]
			The report must explicitly identify \emph{what was executed} (e.g., a reconstructed function call) rather than relying on the test name or on implicit context.

		\item[\textbf{R3: Expected vs.\ actual as the primary comparison.}]
			The expected and actual outcomes must be presented side-by-side (or in a tightly coupled pair) in a way that supports rapid visual comparison.

		\item[\textbf{R4: Separation of essential and auxiliary information.}]
			Framework internals (stack traces, runner frames, reflection calls) must not obscure the essential summary. If included at all, such details should be clearly separated and optionally suppressible.

		\item[\textbf{R5: Applicability across test paradigms.}]
			The same summary structure should apply to both example-based unit tests and property-based tests, even when additional information (e.g., counterexamples) is present.

		\item[\textbf{R6: Structured representation of property-based counterexamples.}]
			If a property fails, the report must clearly indicate (i) the failing input, (ii) any shrunk/minimised counterexample, and (iii) how this input relates to the tested expression and observed mismatch.

		\item[\textbf{R7: Support for contextual history when needed.}]
			For multi-step/stateful scenarios, the reporting layer should be able to attach an execution history (e.g., a sequence of operations) to make the failure reproducible and interpretable.
	\end{description}
\end{justify}


%===============================================================================
\subsection*{Design implications for Assertify/Checkify}
%===============================================================================


\begin{justify}
	The requirements above imply that a reporting layer must (a) gain access to the tested expression at runtime,
	(b) evaluate it to obtain the actual value, (c) retain or compute the expected value, and (d) serialise both
	into a stable output model that is independent of any particular baseline framework.

	\skipM In the scope of this thesis, these implications motivate:

	\skipS\begin{itemize}
		\item \textbf{Expression capture:} obtain a representation of the tested expression that can be displayed to the user (and simplified where appropriate).
		\item \textbf{Value extraction:} evaluate expressions in a controlled way to compute actual values and to render values consistently.
		\item \textbf{Wrapper-based integration:} intercept failures from existing unit testing and property-based testing frameworks and replace or
			augment their output without requiring changes to the surrounding tooling.
		\item \textbf{Unified output model:} define a single report representation that can be produced for different assertion forms and test paradigms.
	\end{itemize}

	\skipS\textbf{Checkpoint (traceability to later chapters):}

	\skipS\begin{itemize}
		\item TODO: Add a small mapping table from issues $\rightarrow$ requirements $\rightarrow$ implementation mechanisms (one row per requirement).
		\item TODO: Decide which requirements are \emph{must-have} vs.\ \emph{nice-to-have} for the first version of Assertify.
		\item TODO: State explicitly which requirements are out of scope (e.g., improving compiler diagnostics beyond test reporting).
	\end{itemize}
\end{justify}


%===============================================================================
\subsection*{Completion Checklist for This Chapter}
%===============================================================================


\begin{justify}
	\begin{itemize}
		\item Add at least one fully worked example for each failure category (compiler error, unit test failure, property-based failure).
		\item Add an explicit ``signal vs.\ noise'' annotation for each example output.
		\item Ensure every derived requirement is supported by at least one observed issue (traceability).
		\item Add short, properly cited connections to prior work on diagnostic usability and novice comprehension (do not overclaim; separate observation from generalisation).
		\item Verify terminology consistency with later chapters (e.g., always use \emph{tested expression / expected / actual} in the same order).
	\end{itemize}
\end{justify}


%===============================================================================
% TODO +++++++++++++++++++++++++++++++++++++++++++++++++++++++++++++++++++++++++
%===============================================================================


% TODO


%===================================================================================================
% EOF ++++++++++++++++++++++++++++++++++++++++++++++++++++++++++++++++++++++++++++++++++++++++++++++
%===================================================================================================